% Application specification shared by the report and standalone fiches.
\ApplicationFicheHeading{app-feelpp-parmmg}
\label{sec:app:specs:app-feelpp-parmmg}

% Generated by scripts/sync_applications.py; edit the CSV snapshot.
\begin{FicheMetadataTable}
id & app-feelpp-parmmg \\
name & Parallel remeshing with ParMmg \\
PC & PC1 - ExaMA \\
work\_package & WP1 \\
application\_type & mini-app \\
purpose & Evaluate Feel++/ParMmg parallel mesh adaptation; case definition and benchmark evidence pending. \\
Method-Algorithm WP1 & mesh adaptation, unstructured mesh \\
status & spec-only \\
Framework & Feel++, MMG/ParMMG \\
parallel\_framework & MPI \\
repo\_url & https://github.com/feelpp/feelpp \\
notes & Separate Feel++ application, distinct from app-freefem-parmmg. Upstream Feel++ supplies the implementation context; application configuration, contributors, verification protocol and benchmark evidence remain to be supplied. \\
\end{FicheMetadataTable}


\subsection{Description of the benchmark}
This application covers Feel++/ParMmg parallel mesh adaptation, separately
from the FreeFEM++/ParMmg demonstrator. Its catalogue status is
\texttt{spec-only}; the case, contributors and benchmark data remain pending.

\href{https://github.com/feelpp/feelpp}{Feel++} provides a ParMmg remesher and
remeshing test drivers. These provide a starting point; application-level
verification and scaling measurements are still to be defined.

\paragraph{Planned contribution.} Identify the intended Feel++ case,
executable and configuration. Confirm whether this is a remeshing kernel or
a coupled solve--adapt workflow.


\paragraph{Planned remeshing contract.}
The owner must select either a remeshing kernel (mesh plus metric to adapted
mesh) or a solve--estimate--adapt loop. For the kernel, supply the spatial
metric tensor convention/units, anisotropy bounds, gradation and target-size
settings. Preserve required physical markers and interfaces; check positive
cell orientation/volume, metric quality and achieved element sizes. Record
before/after cell counts, quality distributions and preparation, adaptation
and redistribution times. For a coupled loop, additionally define error
estimation and field-transfer/conservation checks and compare at fixed error.
A configuration of an existing Feel++ remeshing driver can implement this
contract; a new C++ application is required only if that driver lacks the
selected operation. The FreeFEM Fichera configuration is a separate case.

\subsection{Benchmarking methodology}
\paragraph{Planned contribution.} Specify the mesh and metric, verification
checks, baseline, parallel resource series, measured phases and repetition
protocol. Pin the application configuration and software environment before
reporting measurements.

\subsection{Input/Output Dataset Description}
\subsubsection{Input Data}
\paragraph{Planned contribution.} Supply the input geometry or mesh,
adaptation metric, required markers and remeshing parameters for the chosen case.

\subsubsection{Output Data}
\paragraph{Planned contribution.} Document the adapted-mesh outputs,
verification reports and timing records actually produced by the application.

\subsection{Results summary}
No measured results or benchmark-readiness claim are assigned to this draft.
\paragraph{Planned contribution.} Attach versioned driver/configuration,
validation reports and campaign artefacts when available; keep them planned until then.

% Generated by scripts/feelpp_catalogue.py; edit data/feelpp-mesh-catalogue.json.

\subsection{Application package and shared meshes}
\begin{FicheMetadataTable}
Repository entry & \nolinkurl{apps/app-feelpp-parmmg/} (local review layout in apps-feelpp; publication pending) \\
Delivery model & driver config planned \\
Available material & Feel++ remesher and upstream test driver available; dedicated case not selected. \\
Results status & Feel++ campaign results planned. \\
\hyperref[mesh:feelpp:parmmg-case]{Feel++ ParMmg case} & Geometry and physical units pending case selection. Choose case, metric convention and parameters, protected markers, mesh-quality and interpolation checks. \\
\end{FicheMetadataTable}
\paragraph{Next contribution.} Choose remeshing kernel or solve-adapt workflow; supply input mesh/metric, boundary markers and quality/conservation checks. Do not reuse the FreeFEM Fichera fiche.

